\documentclass[11pt,a4paper]{article}
\usepackage[margin=23mm]{geometry}
\usepackage{amsmath,amssymb,booktabs,array}
\input{glyphtounicode}
\pdfgentounicode=1
\usepackage[hidelinks]{hyperref}
\setlength{\parindent}{0pt}
\setlength{\parskip}{6pt}
\begin{document}
\begin{center}
{\Large K2: an arrival-aware compulsory merge}\\[4pt]
14 October 2026 --- conditional candidate and interface gate, not adopted methods
\end{center}

\section*{1. Fixed geometry, information and the proposed handoff}
Keep total road length $L$ fixed. Candidate T1 reserves a terminal reach
$b\in(0,L)$ and applies homogeneous early preparation on $[0,L_p]$, where
$L_p=L-b$. An initial uniform footprint $[0,\ell]$, $0<\ell\leq L_p$, has
ending, adjacent and other-continuing densities $e_0,a_0,b_0$, with
$a_0=b_0$ and $M=\ell(e_0+a_0+b_0)>0$. The terminal reach is initially empty.
There are no later arrivals or internal lateral sources in that reach.
All quantities and the handoff are fixed across advice actions and comparators.
The geometry and cessation of early transfer at $L_p$ are proposals requiring
adoption, not consequences of approving K2's review direction.

Use one common action $v\in[v_-,v_+]$, $0<v_-<v_+$, and the H1 diagnostic
family $k_c(v)=wk_j/(v+w)$, $C(v)=vk_c(v)$, with $w,k_j>0$. Preserve the H1
initial bounds $0\leq e_0,a_0,b_0\leq k_c(v_+)$. Its homogeneous body evolves as
\[
 r=e_0+a_0=e(t)+a(t),\qquad
 \dot e=-\frac{v}{h}\min\left\{e,\frac{k_c(v)-r+e}{\gamma}\right\},
 \quad a=r-e,\quad b(t)=b_0,
 \tag{1}
\]
with fixed physical interaction scale $h>0$ and early work weight $\gamma\geq1$.
The H1 proof of nonnegative mass, common-speed free flow and lane-density
invariance applies only while downstream admission does not hold the body.
The adaptations in (1) remain as declared in the H1 note.

Let $F_E,F_A,F_B$ be cumulative terminal-entry mass, initially zero. The
preparation front and rear leave at $(L_p-\ell)/v$ and $L_p/v$. During that
interval,
\[
 \dot F_E=ve(t),\qquad \dot F_A=va(t),\qquad \dot F_B=vb_0;
 \tag{2}
\]
outside it all entry rates are zero. Early transfers into $A$ are already in
$F_A$ and are not forced to merge a second time. Route-origin conservation
may be checked by paired early transfers; service needs only the current
completed/uncompleted class.

\section*{2. Physical stock, eligible queue and receiving at the interface}
The terminal reach has free-flow delay $\tau=b/v$, backward delay
$\tau_b=b/w$ and per-lane jam storage $H=k_jb$. Let $Y_j$ be cumulative final
departures for $j=E,A,B$. Extend initial histories as zero for negative time.
The initially empty-link quantities are
\[
 N_j(t)=F_j(t)-Y_j(t),\quad
 P_j(t)=F_j(t-\tau),\quad
 q_j(t)=P_j(t)-Y_j(t).
 \tag{3}
\]
$N_j$ is physical stock; $P_j$ is mass that could have reached the lane end
without delay; $q_j$ is eligible but unserved mass. These are different
objects. A virtual eligibility time is not the physical moment a traveler
joins the back of a spatial queue.

For $t>\tau_b$, the LTM vacancy is
\[
 V_j(t)=Y_j(t-\tau_b)+H-F_j(t).
 \tag{4}
\]
Positive vacancy permits upstream receiving at capacity. At $0<t\leq\tau_b$,
the empty-link vacancy is $k_jwt-F_j(t)$; capacity-bounded input makes it
positive. The input rate must also respect $C(v)$. A conservative sufficient
domain for all lanes and all actions is
\[
 M<H=k_jb,\qquad 0\leq\dot F_j(t)\leq C(v).
 \tag{5}
\]
Indeed $F_j\leq M<H$ and $Y_j\geq0$ ensure (4) is strictly positive after
the backward delay. Before it, $F_j(t)\leq C(v)t<k_jwt$. Thus terminal
admission never restricts H1. This is a finite, deliberately conservative
no-spillback domain, not an infinite buffer and not a test on $q_E$ alone.

Eligible demand at the ending-lane exit is zero before $\tau$. Afterwards it
is its new eligible-arrival rate when $q_E=0$ and its link capacity $C(v)$
when $q_E>0$. This is the empty-link LTM queue/demand relation. The formulas
are a restricted source-compatible reference; no assertion is made that
the weighted node below equals Jin's general invariant junction model.

\section*{3. One queue, own-lane priority and finite forced service}
Let $R>0$ be declared constant available receiving-work rate for the adjacent
exit and $R_B>0$ the other exit's capacity. Ordinary mass uses one work unit;
unfinished ending mass uses $\alpha\geq1$. The late factor $\alpha$ is separate
from $\gamma$. Select the elementary nominal domain
\[
 R>v_+r,\qquad R_B\geq v_+b_0.
 \tag{6}
\]
This is stronger than merely fitting current continuing demand: the whole
ordinary $E+A$ load fits $R$. Any forced queue arises from its extra receiving
work, rather than an unacknowledged ordinary load excess. Source sending is
respected because each lane's input is on the H1 free-flow branch.

Own-lane-first service then gives $Y_A=P_A$, $Y_B=P_B$ and
$u_A=\dot P_A$. The remaining ending-lane service limit is
\[
 c_E(t;v)=\min\left\{C(v),\frac{R-u_A(t)}{\alpha}\right\}>0.
 \tag{7}
\]
The terminal flow and queue balance are
\[
 g_E(t)=\dot Y_E(t)=
 \begin{cases}
 c_E(t),&q_E(t)>0,\\
 \min\{\dot P_E(t),c_E(t)\},&q_E(t)=0,
 \end{cases}
 \qquad \dot q_E=\dot P_E-g_E.
 \tag{8}
\]
Hence $q_E\geq0$, $Y_E\leq P_E$, $g_E$ respects ending sending, and
\[
 u_A+\alpha g_E\leq R,\qquad
 u_A+g_E\leq R.
 \tag{9}
\]
Transferred/departed mass is $g_E\,dt$, not $\alpha g_E\,dt$. Forced mass
has no forward escape branch. Failed requests remain in (3), and the source
link represents their slower congested motion implicitly. Advice is still
the free-flow parameter; queued terminal traffic is not asserted to travel
at that speed.

For the work-conserving diagnostic below, additionally require
\[
 R/\alpha\leq C(v_-).
 \tag{10}
\]
Since $C(v)$ increases with $v$, the longitudinal cap in (7) cannot bind in
that subdomain. It is retained in the general candidate instead of hidden.
Conditions (5), (6) and (10) are restrictions to qualify over the whole action
range, not post-action receiving curves.

The entire finite cohort has become eligible by $t_*=L/v$. Thereafter own
arrivals are zero and positive service is $\min\{C(v),R/\alpha\}$. Thus
\[
 T_{\rm clear}=t_*+
 \frac{q_E(t_*)}{\min\{C(v),R/\alpha\}},
 \tag{11}
\]
with the no-queue case interpreted as $T_{\rm clear}=t_*$. A positive native
continuing density supplies the last physical tail; otherwise (11) is a
finite upper bound. No unfinished mass is discarded at a horizon.

\textbf{Source versus adaptation.} Local own-lane-first and gap-weighted
residual receiving are Carey ingredients. The compulsory boundary extension,
constant terminal work factor/rate, restricted priority domain, final reach
and lack of lateral sources there are explicit adaptations. Cooperative or
proportional service is a different terminal rule, not a harmless replacement
of (8); it can introduce an $A$ queue. No field calibration or zipper law is
claimed.

\section*{4. A joined domain with positive waiting and no upstream restriction}
All the nominal restrictions can hold together. Choose positive sufficiently
sparse initial densities so (5), (6) and (10) have room. At an admitted speed
$v_0$, ending density at the first terminal arrival is positive because its
finite H1 exposure has not annihilated all mass. Let those eligible densities
be $e_f,a_f$, with $e_f+a_f=r$.

Positive initial queue growth occurs whenever
\[
 v_0r<R<v_0(a_f+\alpha e_f)
       =v_0[r+(\alpha-1)e_f].
 \tag{12}
\]
For $v_0=v_+$ and $\alpha>1$, this is a nonempty interval above the ordinary
load bound. Sparse traffic can also make $R/\alpha<C(v_-)$ while $M<k_jb$.
These are compatibility/existence inequalities, not fitted parameters or
an interior-optimum certificate. Receiving at the upstream handoff remains
slack under (5), even while the compulsory node has a queue under (12).

If (5) or the actual vacancy criterion fails, H1's prescribed entry clock
is no longer trusted. If (6) fails, the one-queue restriction may no longer
hold. Those failures reopen the smallest interface/domain contract; neither
clipping incoming mass nor adding an undocumented second queue is permitted.

\section*{5. Complete TST, conservation and an arrival-pulse reference}
All cohort mass is initially present at decision time. Total unexited mass
and complete TST are
\[
 N_{\rm unexited}(t)=M-Y_E(t)-Y_A(t)-Y_B(t),\qquad
 \mathrm{TST}=\int_0^{T_{\rm clear}}N_{\rm unexited}(t)\,dt.
 \tag{13}
\]
Paired early transfer conserves upstream mass, (3) conserves each terminal
stream's input/output stock, and node discharge removes exactly the accepted
mass. Gap work is never multiplied into population counts.

Because every initial position has free-flow eligibility time $(L-x)/v$
and only $E$ is delayed in this domain,
\[
 \mathrm{TST}=M\frac{L-\ell/2}{v}
                 +\int_0^{T_{\rm clear}}q_E(t)\,dt.
 \tag{14}
\]
This is an identity for the admitted full trajectory. It counts physical
travel through both regions and all excess terminal delay. It does not place
a positive spatial queue after the road or claim that (3)'s eligible queue
equals physical occupancy.

For an isolated node reference, let masses $m_A=m_E=4$ arrive uniformly over
duration $d$, set $R=4$, $\alpha=2$, and let sending be nonbinding. The
reference clock starts at the first eligible arrival, not a newly chosen
H1 decision boundary. During arrivals
\[
 u_A=m_A/d,\quad c_E=(R-m_A/d)/\alpha,\quad
 q_{\rm peak}=\frac{m_A+\alpha m_E-Rd}{\alpha}
 \tag{15}
\]
when the numerator is positive. If it is nonpositive, no queue forms.
After arrivals the drain duration is $\alpha q_{\rm peak}/R$. The two
triangular queue areas sum to
\[
 \Delta\mathrm{TST}=\frac{q_{\rm peak}d}{2}
              +\frac{\alpha q_{\rm peak}^2}{2R}.
 \tag{16}
\]
For $d=2$, $q_{\rm peak}=2$, drain duration is one, delay is three and the
eligible-arrival-time sum $(m_A+m_E)d/2$ is eight: complete-time sum is eleven.
For $d=4$, the delay is zero and the complete-time sum is sixteen. These are
closed-form hand references, not traffic rollouts, calibrated actions, or a
reachable H1 comparison. Any whole-road prefix must still be included.

\section*{6. Priority reference: equal final work is not equal TST}
Consider a different, already waiting batch with $m_A=m_E=2$, $R=2$ and
$\alpha=2$; all sending limits are slack. This comparison deliberately allows
an $A$ backlog and is outside the nominal one-queue input domain. Total work
is six, so both following work-conserving allocations finish at time three.

Own-lane-first clears $A$ at rate two over $[0,1]$, then $E$ at rate one over
$[1,3]$. The remaining population integrates to
\[
 \int_0^1(4-2t)\,dt+\int_1^3(3-t)\,dt=3+2=5.
 \tag{17}
\]
Proportional work allocation maintaining the batch fractions serves both
mass classes at rate $2/3$ over $[0,3]$, and
\[
 \int_0^3 4(1-t/3)\,dt=6.
 \tag{18}
\]
Thus S13's proportional batch is an actual conditional rule, not merely the
name of a component. Its mean is not the mean of an own-lane-first queue.
This reference neither calibrates priorities nor declares the cooperative
alternative invalid. A full arrival-aware rival must specify its available
demand, priority and own queue explicitly.

\section*{7. Fixed-preparation speed diagnostic: fastest can remain best}
Hold completed preparation outcomes fixed across actions: each initial
position carries the same terminal service class for every $v$. This is a
deletion/reference case, not an assumption for the actual H1 sensitivity.
For $v_f>v_s$, all free-flow eligibility times are earlier, so for all $t$
\[
 P_{A,f}(t)\geq P_{A,s}(t),\quad
 P_{E,f}(t)\geq P_{E,s}(t),\quad
 P_{B,f}(t)\geq P_{B,s}(t).
 \tag{19}
\]
Assume the same fixed $R,\alpha$ and domain (5), (6), (10). Put
\[
 W(t)=P_A(t)+\alpha P_E(t),\qquad
 Z(t)=Y_A(t)+\alpha Y_E(t).
\]
The weighted server is work-conserving: whenever weighted backlog is positive,
it serves work at $R$. With zero initial eligible mass its cumulative output is
\[
 Z(t)=\inf_{0\leq u\leq t}\{W(u)+R(t-u)\}.
 \tag{20}
\]
This exact fluid reference follows from conservation and reflected queue
service. Earlier cumulative input (19) implies $Z_f(t)\geq Z_s(t)$. Since
$Y_A=P_A$, total completed mass is
\[
 Y_A+Y_E+Y_B=\frac{Z}{\alpha}
                   +\left(1-\frac1\alpha\right)P_A+P_B.
 \tag{21}
\]
Every coefficient is nonnegative. Hence faster advice has at least as much
completed mass at every time, and its integral (13) is no greater:
\[
 \mathrm{TST}(v_f)\leq\mathrm{TST}(v_s).
 \tag{22}
\]
Forced queueing alone therefore does not guarantee an interior minimum.
Changed H1 preparation changes the input work/classes and invalidates the
fixed-preparation premise of this proof. An actual interior optimum, if any,
needs a full sensitivity/evidence pass and is not imposed as an acceptance
criterion for this server.

As a separate deletion, setting $\alpha=1$ under the ordinary-load bound (6)
makes all arrivals fit the receiver and their own sending, so no ending queue
forms. Complete time is then only the free-flow term in (14). This deletion
does not claim aggregate equivalence in a broader domain that violates (6)
or adds other receiver asymmetry, priorities or capacity-drop mechanisms.

\section*{8. Scientific disposition and references}
Initial symmetry and an empty terminal mean $M$ plus $E_0$ fixes
$e_0=E_0/\ell$, $a_0=b_0=(M-E_0)/(2\ell)$ and all candidate initial states.
The strong aggregate-plus-ending-statistic rival can predict the same histories
and queue. It is given all fixed parameters and the same observation contract.
No future receiving oracle or extra initial lane observable is introduced.

The physical T1 construction and explicitly virtual T2 server are different
methodological interpretations. T1 needs adoption of the terminal reach,
early-transfer cutoff, proposed priority and joined domain; T2 needs acceptance
of a deliberate no-feedback approximation. Neither exact law is adopted here.
No applied joint calibration, reachable new information witness, AD/finite-
difference evidence or independent usefulness result is claimed. The K2 pass
stops at the original B2 interface gate; the modeling-chain stage remains open.

\begin{itemize}
\item Carey, Balijepalli and Watling (2013), Sections 5 / R1, R2 and Algorithms
A1--A2: local receiving/priority ingredients.
\href{https://eprints.whiterose.ac.uk/id/eprint/77315/7/Carey\%20et\%20al\%202013_with_coversheet.pdf}{Institutional manuscript}.
\item Jin (2015), Sections 3.1--3.3, equations (23), (28)--(30): empty-link
queue, vacancy and boundary counts.
\href{https://arxiv.org/pdf/1405.7080}{Primary preprint}.
\item Yperman, Logghe and Immers (2005), Sections 3--4: link/node decomposition
and merge allocation.
\href{https://www.mech.kuleuven.be/cib/verkeer/dwn/pub/P2005C.pdf}{Original LTM paper}.
\end{itemize}
All terminal geometry, no-lateral-source, weighted priority and action-domain
adaptations are declared explicitly. The accepted S13 certificate remains
unchanged.
\end{document}
